\documentclass[11pt]{article}

\usepackage[a4paper,margin=1.05in]{geometry}
\usepackage{amsmath,amssymb,amsthm,mathtools}
\usepackage{enumitem}
\usepackage{hyperref}
\usepackage{mathrsfs}

\hypersetup{
  colorlinks=true,
  linkcolor=blue,
  citecolor=blue,
  urlcolor=blue
}

\numberwithin{equation}{section}

\newtheorem{theorem}{Theorem}[section]
\newtheorem{lemma}[theorem]{Lemma}
\newtheorem{proposition}[theorem]{Proposition}
\newtheorem{corollary}[theorem]{Corollary}
\newtheorem{assumption}[theorem]{Assumption}
\newtheorem{definition}[theorem]{Definition}
\newtheorem{remark}[theorem]{Remark}

\newcommand{\R}{\mathbb R}
\newcommand{\eps}{\varepsilon}
\newcommand{\loc}{\mathrm{loc}}
\newcommand{\CKN}{\mathrm{CKN}}
\newcommand{\RNSE}{\mathrm{RNSE}}
\newcommand{\NS}{\mathrm{NS}}
\newcommand{\cE}{\mathcal E}
\newcommand{\cS}{\mathcal S}
\newcommand{\cA}{\mathcal A}
\newcommand{\cM}{\mathcal M}
\newcommand{\NoMulti}{\mathrm{NoMulti}_{R4}}

\title{Discharging the \(R4\) Branch from a Finite No-Multibubble Theorem}
\author{}
\date{}

\begin{document}
\maketitle

\begin{abstract}
This note isolates the final logical step in the \(R4\) branch.  We prove, in
standard PDE language, that the generic \(R4\) residual class is empty once
finite retained active multibubble terminal states are excluded.  We also state
the precise finite no-multibubble theorem and show how it may be reused from a
terminal nonlinear profile-decoupling theorem.  The proof uses terminal
state-space exhaustion, inactive-stratum discharge, recurrent tail rigidity,
atomic sequence-\(L^3\) completeness, mildness inheritance, and the
Albritton--Barker Liouville theorem for mild ancient Navier--Stokes solutions.
\end{abstract}

\tableofcontents

\section{Renormalized setting}

We work with the backward self-similar renormalized Navier--Stokes equation
\begin{equation}
  \partial_\tau U-\Delta U+\frac12 y\cdot\nabla U+\frac12 U
  +(U\cdot\nabla)U+\nabla\Pi=0,
  \qquad
  \nabla\cdot U=0
  \label{eq:rnse}
\end{equation}
on \(\R^3_y\times\R_\tau\).  The associated physical variables are
\begin{equation}
  u(x,t)=(-t)^{-1/2}
  U\!\left(\frac{x}{\sqrt{-t}},-\log(-t)\right),
  \qquad t<0,
  \label{eq:pullback-u}
\end{equation}
and
\begin{equation}
  p(x,t)=(-t)^{-1}
  \Pi\!\left(\frac{x}{\sqrt{-t}},-\log(-t)\right).
  \label{eq:pullback-p}
\end{equation}
A direct computation gives that \((u,p)\) solves the physical Navier--Stokes
equations
\begin{equation}
  \partial_tu-\Delta u+(u\cdot\nabla)u+\nabla p=0,
  \qquad
  \nabla\cdot u=0.
  \label{eq:physical-ns}
\end{equation}

For a parabolic cylinder
\[
  Q_r(z_0):=B_r(y_0)\times(\tau_0-r^2,\tau_0),
  \qquad z_0=(y_0,\tau_0),
\]
define the local Caffarelli--Kohn--Nirenberg quantity
\begin{equation}
  \cE_r(U,\Pi;z_0)
  :=
  \iint_{Q_r(z_0)}
  \left(
    |U|^3+\left|\Pi-c_{r,z_0}(\tau)\right|^{3/2}
  \right)\,dy\,d\tau,
  \label{eq:local-ckn}
\end{equation}
where \(c_{r,z_0}(\tau)\) is an admissible local pressure gauge, for instance
the average of \(\Pi(\cdot,\tau)\) over \(B_r(y_0)\).  We write \(Q_r=Q_r(0,0)\).

\section{The \(R4\) class and terminal decomposition}

The class \(R4\) is the residual Type I branch left after all previously
identified branches have been removed.  Its elements are bounded ancient
Seregin-limit profiles satisfying the Type I \(L^\infty\) bound, local
suitability, persistent compact CKN density, and none of the previously closed
small, stationary \(L^3\), tight, fast-decay, structured, axisymmetric,
controlled-swirl, R1, R2, or R3 alternatives.

\begin{definition}[R4 profile]
An \(R4\) profile is a smooth bounded ancient solution \((V,P)\) of
\eqref{eq:rnse} such that:
\begin{enumerate}[label=(\roman*)]
  \item \(\|V\|_{L^\infty(\R^3\times\R)}\le M\);
  \item \(V\) is suitable on every compact cylinder;
  \item \(V\) carries persistent compact CKN density, i.e. along the terminal
  extraction under consideration there is a compact cylinder on which the
  local CKN quantity is bounded below by a fixed positive number;
  \item \(V\) is not in any of the previously closed branches.
\end{enumerate}
\end{definition}

We assume the following terminal decomposition input.

\begin{assumption}[Terminal state-space exhaustion]
\label{ass:terminal-exhaustion}
Every terminal extraction of a tail-normalized \(R4\) profile belongs to exactly
one of the following alternatives:
\[
  \text{core},\quad
  \text{scattering},\quad
  \text{exterior},\quad
  \text{radiative},\quad
  \text{rough},\quad
  \text{finite multibubble}.
\]
Furthermore, the scattering, exterior, radiative, and rough alternatives do not
retain the compact active CKN density of an \(R4\) profile.
\end{assumption}

\begin{assumption}[Tail reduction and affine normalization]
\label{ass:tail-normalization}
The tail-hull analysis has been performed.  Nonconstant recurrent tail cores
are discharged by recurrent tail rigidity, and constant tails are removed by
the affine renormalized Galilean symmetry.  Thus any remaining \(R4\) candidate
may be represented as a tail-zero profile.
\end{assumption}

\section{Retained active profiles and finite multibubbles}

Fix a small positive threshold \(\eps_*>0\).

\begin{definition}[retained active terminal profile]
A terminal profile \((U,\Pi)\) is called retained active if
\begin{equation}
  \cE_1(U,\Pi;0)\ge \eps_*.
  \label{eq:active-profile}
\end{equation}
\end{definition}

\begin{definition}[finite retained active multibubble]
A finite retained active multibubble terminal state is a finite family of
terminal cameras
\[
  g_{1,n},\dots,g_{J,n},
  \qquad J\ge2,
\]
such that:
\begin{enumerate}[label=(\roman*)]
  \item the cameras are mutually separated:
  \[
    g_{i,n}^{-1}g_{j,n}\to\infty
    \qquad (i\ne j);
  \]
  \item after passing to a subsequence,
  \[
    g_{j,n}\cdot V\to U_j
  \]
  locally smoothly for velocity and locally in \(L^{3/2}\) for pressure modulo
  admissible gauges;
  \item each terminal profile is retained active:
  \[
    \cE_1(U_j,\Pi_j;0)\ge\eps_*,
    \qquad j=1,\dots,J.
  \]
\end{enumerate}
The packet is called maximal if no further separated retained active terminal
camera may be added to it.
\end{definition}

\begin{theorem}[Finite no-multibubble theorem]
\label{thm:no-multi}
No \(R4\) profile admits a finite retained active multibubble terminal state.
\end{theorem}

Theorem~\ref{thm:no-multi} is the final multibubble statement.  In
Section~\ref{sec:reuse}, we explain how it can be reused from terminal
nonlinear profile decoupling.

\section{Atomic terminal profiles}

\begin{definition}[tail descendant]
Let \((U,\Pi)\) be a terminal profile.  A retained active tail descendant of
\(U\) is a terminal limit \(W\) obtained from cameras escaping to infinity
relative to \(U\), such that
\[
  \cE_1(W,\Pi_W;0)\ge\eps_*.
\]
\end{definition}

\begin{definition}[radiative and rough descendants]
A radiative descendant of \(U\) is a sequence
\[
  \tau_n\to-\infty,\qquad R_n\to\infty,
\]
and a number \(\eta>0\), such that
\[
  \int_{|y|>R_n}|U(y,\tau_n)|^3\,dy\ge\eta,
  \label{eq:radiative}
\]
while every unit-scale exterior camera is below the retained-active threshold,
after excluding rough cameras.

A rough descendant is a tail camera for which the local compactness,
suitability, or pressure control needed to extract a terminal profile fails.
\end{definition}

\begin{definition}[atomic profile]
A retained active terminal profile is atomic if it has no retained active tail
descendant, no radiative descendant, and no rough descendant.
\end{definition}

\section{One active descendant gives a two-camera multibubble}

\begin{lemma}
\label{lem:desc-to-multi}
Let \(U\) be a retained active terminal profile.  If \(U\) has a retained
active tail descendant, then \(U\) belongs to the finite retained active
multibubble stratum.
\end{lemma}

\begin{proof}
The origin camera of \(U\) carries positive local CKN mass:
\[
  \cE_1(U,\Pi;0)\ge\eps_*.
\]
Let \(W\) be a retained active tail descendant.  By definition there are
cameras \(h_n\to\infty\), relative to \(U\), such that
\[
  h_n\cdot U\to W
\]
in the terminal topology and
\[
  \cE_1(W,\Pi_W;0)\ge\eps_*.
\]
By local convergence of the velocity and pressure, after gauges,
\[
  \cE_1(h_n\cdot U,\Pi_{h_n};0)\ge \frac12\eps_*
\]
for all sufficiently large \(n\).  Thus, for large \(n\), the same terminal
object has two separated retained active cameras: the origin camera and the
camera \(h_n\).  This is a finite two-camera retained active multibubble.
\end{proof}

\begin{proposition}
\label{prop:no-multi-implies-atomic}
Assume Theorem~\ref{thm:no-multi} and
Assumption~\ref{ass:terminal-exhaustion}.  Then every singleton retained active
terminal profile arising from an \(R4\) candidate is atomic.
\end{proposition}

\begin{proof}
Let \(U\) be a singleton retained active terminal profile.  If \(U\) had a
retained active tail descendant, Lemma~\ref{lem:desc-to-multi} would place it
in the finite multibubble stratum, contradicting Theorem~\ref{thm:no-multi}.

If \(U\) had a radiative or rough descendant, then the corresponding terminal
extraction would lie in the radiative or rough alternative.  These alternatives
are excluded as carriers of the persistent compact CKN density by
Assumption~\ref{ass:terminal-exhaustion}.  Hence no such descendant exists, and
\(U\) is atomic.
\end{proof}

\section{Atomic sequence-\(L^3\) completeness}

\begin{lemma}
\label{lem:atomic-L3}
Let \(U\) be a bounded smooth ancient solution of \eqref{eq:rnse}.  If \(U\) is
retained active and atomic, then there exists a sequence
\[
  \tau_k\to-\infty
\]
such that
\[
  \sup_k\|U(\cdot,\tau_k)\|_{L^3(\R^3)}<\infty.
\]
\end{lemma}

\begin{proof}
We argue by contrapositive.  Suppose no such sequence exists.  Then
\begin{equation}
  \liminf_{\tau\to-\infty}
  \|U(\cdot,\tau)\|_{L^3(\R^3)}
  =\infty.
  \label{eq:no-seq-L3}
\end{equation}
Choose \(\tau_n\to-\infty\) so that
\[
  \int_{\R^3}|U(y,\tau_n)|^3\,dy\to\infty.
\]
Since \(U\) is bounded, choose \(R_n\to\infty\) sufficiently slowly so that
\begin{equation}
  \int_{|y|>R_n}|U(y,\tau_n)|^3\,dy\ge 1.
  \label{eq:tail-mass}
\end{equation}
Set
\[
  m_n:=
  \sup_{|x|>R_n}
  \int_{B_1(x)}|U(y,\tau_n)|^3\,dy.
\]

If \(\limsup_n m_n\ge\eps_*\), then for a subsequence there exist
\(|x_n|>R_n\) such that
\[
  \int_{B_1(x_n)}|U(y,\tau_n)|^3\,dy\ge\eps_*.
\]
After recentering at \((x_n,\tau_n)\), either terminal compactness fails,
which gives a rough descendant, or a terminal subsequence converges to a
retained active tail descendant.  Both alternatives contradict atomicity.

It remains to consider the case \(m_n<\eps_*\) for all sufficiently large
\(n\).  Then \eqref{eq:tail-mass} says that exterior \(L^3\)-mass persists,
while every unit-scale exterior camera is below the retained-active threshold.
This is precisely a radiative descendant, contradicting atomicity.

Thus \eqref{eq:no-seq-L3} is impossible.
\end{proof}

\section{Albritton--Barker endpoint}

\begin{assumption}[mildness inheritance]
\label{ass:mildness}
If \(U\) is a retained active terminal profile, then the physical pullback
\eqref{eq:pullback-u} is a mild ancient solution of the physical
Navier--Stokes equations, after any finite time shift.
\end{assumption}

\begin{theorem}[Albritton--Barker]
\label{thm:AB}
Let \(u\) be a mild ancient solution of the three-dimensional Navier--Stokes
equations.  If there exists a sequence \(t_k\downarrow-\infty\) such that
\[
  \sup_k\|u(\cdot,t_k)\|_{L^3(\R^3)}<\infty,
\]
then \(u\equiv0\).
\end{theorem}

\begin{lemma}
\label{lem:atomic-vanish}
Assume Assumption~\ref{ass:mildness}.  No retained active atomic terminal
profile exists.
\end{lemma}

\begin{proof}
Let \(U\) be retained active and atomic.  By Lemma~\ref{lem:atomic-L3}, choose
\(\tau_k\to-\infty\) such that
\[
  \sup_k\|U(\cdot,\tau_k)\|_{L^3}<\infty.
\]
Let \(u\) be the physical pullback \eqref{eq:pullback-u}, and set
\[
  t_k=-e^{-\tau_k}.
\]
Then \(t_k\downarrow-\infty\), and by critical scaling,
\[
  \|u(\cdot,t_k)\|_{L^3(\R^3)}
  =
  \|U(\cdot,\tau_k)\|_{L^3(\R^3)}.
\]
By Assumption~\ref{ass:mildness}, \(u\) is mild ancient.  Theorem~\ref{thm:AB}
therefore gives \(u\equiv0\), hence \(U\equiv0\).  This contradicts retained
activity, since \(\cE_1(U,\Pi;0)\ge\eps_*>0\).
\end{proof}

\section{R4 closure from no multibubbles}

\begin{theorem}[R4 closure from finite no-multibubble]
\label{thm:R4-closure}
Assume:
\begin{enumerate}[label=(\roman*)]
  \item tail reduction and affine normalization,
  Assumption~\ref{ass:tail-normalization};
  \item terminal state-space exhaustion and inactive-stratum discharge,
  Assumption~\ref{ass:terminal-exhaustion};
  \item no finite retained active multibubbles, Theorem~\ref{thm:no-multi};
  \item mildness inheritance, Assumption~\ref{ass:mildness};
  \item the Albritton--Barker Liouville theorem, Theorem~\ref{thm:AB}.
\end{enumerate}
Then the \(R4\) class is empty.
\end{theorem}

\begin{proof}
Suppose, for contradiction, that \(V\in R4\).  By
Assumption~\ref{ass:tail-normalization}, we reduce to a tail-zero representative.
By definition of \(R4\), persistent compact CKN density remains.

Apply terminal state-space exhaustion.  By
Assumption~\ref{ass:terminal-exhaustion}, the persistent compact CKN density
cannot be carried by scattering, exterior, radiative, or rough alternatives.
Hence it is carried either by a finite retained active multibubble terminal
state or by at least one singleton retained active terminal profile.

The first case is impossible by Theorem~\ref{thm:no-multi}.  In the second
case, let \(U\) be such a singleton retained active terminal profile.  By
Proposition~\ref{prop:no-multi-implies-atomic}, \(U\) is atomic.  By
Lemma~\ref{lem:atomic-vanish}, \(U\) cannot exist.  This contradiction proves
that no \(R4\) profile exists.
\end{proof}

\section{Reusing a no-multibubble theorem}
\label{sec:reuse}

The proof above isolates exactly how a no-multibubble theorem is reused in the
\(R4\) branch.  The input required by Theorem~\ref{thm:R4-closure} is simply
Theorem~\ref{thm:no-multi}.  One may prove Theorem~\ref{thm:no-multi} by any
method.  A standard route is through terminal nonlinear profile decoupling.

\begin{theorem}[Terminal nonlinear profile decoupling]
\label{thm:terminal-decoupling}
Let \((V_n,P_n)\) be a terminal sequence with a maximal finite retained active
packet \(g_{1,n},\ldots,g_{J,n}\).  Then, after passing to a subsequence, there
are mild ancient profiles \((U_j,\Pi_j)\), \(j=1,\ldots,J\), such that:
\begin{enumerate}[label=(\roman*)]
  \item \(g_{j,n}\cdot V_n\to U_j\) in \(C^\infty_{\loc}\);
  \item \(g_{j,n}\cdot P_n-c_{j,n}(\tau)\to\Pi_j\) locally in \(L^{3/2}\);
  \item in the \(j\)-th camera, every other profile escapes every compact set;
  \item the pressure generated by distant profiles is harmonic on compact
  cylinders and is absorbed into the local pressure gauge;
  \item nonlinear cross-interactions between distinct profiles vanish on
  compact terminal windows;
  \item after removing the finite packet, no hidden retained active camera
  remains.
\end{enumerate}
\end{theorem}

Theorem~\ref{thm:terminal-decoupling} is the usual terminal no-splitting
statement: separated active profiles behave independently in terminal cameras,
and any further retained active camera contradicts maximality.

\begin{proposition}[Terminal decoupling implies no finite multibubbles]
\label{prop:decoupling-to-no-multi}
Assume Theorem~\ref{thm:terminal-decoupling}, discharge of radiative and rough
terminal alternatives, recurrent tail rigidity, mildness inheritance, and the
Albritton--Barker theorem.  Then Theorem~\ref{thm:no-multi} holds.
\end{proposition}

\begin{proof}
Suppose a finite retained active multibubble packet exists, and let
\(U_1,\ldots,U_J\) be a maximal packet.

If some \(U_j\) is atomic, then Lemma~\ref{lem:atomic-vanish} gives a
contradiction.

Assume therefore that no \(U_j\) is atomic.  Since radiative and rough
alternatives are discharged, every \(U_j\) has a retained active tail
descendant.  By Theorem~\ref{thm:terminal-decoupling}, each such descendant
lifts back to a retained active terminal camera of the original sequence.  By
maximality of the packet, the lifted camera must be equivalent to one of the
existing \(g_{i,n}\).  Thus the finite set of indices
\(\{1,\ldots,J\}\) carries a directed graph in which each vertex has an
outgoing edge.

A finite directed graph with every vertex having an outgoing edge contains a
directed cycle.  This cycle produces a compact recurrent tail set.  By
recurrent tail rigidity, that recurrent tail set consists only of affine
constant profiles.  After tail-zero normalization, the recurrent set is
trivial.  This contradicts retained activity of the profiles in the cycle.

Therefore the assumption that a finite retained active multibubble exists is
impossible.
\end{proof}

Combining Proposition~\ref{prop:decoupling-to-no-multi} with
Theorem~\ref{thm:R4-closure} gives the final reusable form:
\[
  \boxed{
  \text{terminal nonlinear profile decoupling}
  \Longrightarrow
  \text{no finite retained multibubbles}
  \Longrightarrow
  R4=\varnothing.
  }
\]

\section{Summary of the closure chain}

The complete logical chain is:
\[
\begin{aligned}
  &\text{terminal nonlinear profile decoupling} \\
  &\qquad\Longrightarrow
  \text{no finite retained active multibubbles in }R4 \\
  &\qquad\Longrightarrow
  \text{all singleton retained active carriers are atomic} \\
  &\qquad\Longrightarrow
  \text{backward sequence-}L^3\text{ completeness} \\
  &\qquad\Longrightarrow
  \text{Albritton--Barker Liouville contradiction} \\
  &\qquad\Longrightarrow
  R4=\varnothing.
\end{aligned}
\]

\section*{References}

\begin{thebibliography}{9}

\bibitem{AB}
D.~Albritton and T.~Barker,
\emph{On local Type I singularities of the Navier--Stokes equations and
Liouville theorems},
arXiv:1811.00502.

\bibitem{CKN}
L.~Caffarelli, R.~Kohn, and L.~Nirenberg,
\emph{Partial regularity of suitable weak solutions of the Navier--Stokes equations},
Comm. Pure Appl. Math. 35 (1982), 771--831.

\bibitem{Kato}
T.~Kato,
\emph{Strong \(L^p\)-solutions of the Navier--Stokes equation in \(\R^m\), with applications to weak solutions},
Math. Z. 187 (1984), 471--480.

\end{thebibliography}

\end{document}
